%% toc: P3 39 39.6
\section{Proofs for least and partial fixed-point logics}
\label{sec:annex-p03-w03-p03-059-fixed-point-logics}

Throughout, \(\overline X\) is the complement of \(X\) in the ambient set
\(B\), and the notation is that of
\cref{sec:p03-w03-p03-059-fixed-point-logics}.

For the syntactic arguments below, bound first-order and relation variables
are renamed away from free variables and from one another whenever needed.
For a relation binder \([\mathbf{op}\,P\bar y\,.\,\theta](\bar t)\), choose
a fresh symbol \(P'\) of the same arity and rename precisely the
occurrences of \(P\) bound by that operator, stopping at a nested binder
for \(P\). For every interpretation \(Z\), structural induction identifies
the old body with the renamed body when \(P'\) is interpreted by \(Z\).
Thus the two bodies define the same operator, with the same fixed point;
freshness prevents capture and preserves polarity. The same argument,
component by component, applies to a simultaneous fixed-point binder.
In particular, a relation being substituted can be assumed not to be bound
inside the formula, and substitutions concern only its free occurrences.
First-order renaming has the analogous assignment-based justification.

\begin{annexitem}{annex:p03-w03-p03-059-knaster-tarski}{thm:p03-w03-p03-059-knaster-tarski}{Knaster--Tarski for power-set lattices}
\begin{proof}
\emph{Least fixed point.} Let \(S\defeq\{X\subseteq B:F(X)\subseteq X\}\)
and \(L\defeq\bigcap S\). The set \(S\) is nonempty because
\(F(B)\subseteq B\), so \(L\subseteq B\) is well defined.

First, \(F(L)\subseteq L\). Let \(X\in S\). Then \(L\subseteq X\), so
monotonicity gives \(F(L)\subseteq F(X)\subseteq X\). Since \(X\in S\) was
arbitrary, \(F(L)\subseteq\bigcap S=L\).

Second, \(L\subseteq F(L)\). Applying \(F\) to \(F(L)\subseteq L\) gives
\(F(F(L))\subseteq F(L)\), so \(F(L)\in S\) and therefore
\(L=\bigcap S\subseteq F(L)\).

Hence \(F(L)=L\). Every fixed point \(Y\) satisfies
\(F(Y)\subseteq Y\), so \(Y\in S\) and \(L\subseteq Y\). Thus \(L\) is the
least fixed point, and \(L=\bigcap S\). Since \(L\) is itself a fixed
point and is contained in every fixed point,
\(\bigcap\{X:F(X)=X\}=L\) as well.

\emph{Greatest fixed point.} Let \(T\defeq\{X\subseteq B:X\subseteq F(X)\}\)
and \(G\defeq\bigcup T\); \(T\ni\varnothing\) is nonempty. If \(X\in T\),
then \(X\subseteq G\), so \(X\subseteq F(X)\subseteq F(G)\); taking the
union over \(T\) gives \(G\subseteq F(G)\). Applying \(F\) to this
inclusion gives \(F(G)\subseteq F(F(G))\), so \(F(G)\in T\) and
\(F(G)\subseteq G\). Hence \(F(G)=G\). Every fixed point lies in \(T\) and
so is contained in \(G\). Therefore \(G\) is the greatest fixed point and
equals both \(\bigcup T\) and \(\bigcup\{X:F(X)=X\}\).
\end{proof}
\end{annexitem}

\begin{annexitem}{annex:p03-w03-p03-059-closure-ordinal}{thm:p03-w03-p03-059-closure-ordinal}{The stages of an inductive operator stop}
\begin{proof}
\emph{Some stage repeats.} Let \(\kappa\) be the least cardinal greater than
\(|B|\), regarded as an initial ordinal. Suppose that
\(X^\alpha\neq X^{\alpha+1}\) for every \(\alpha<\kappa\). Since \(F\) is
inductive, \(X^\alpha\subsetneq X^{\alpha+1}\), and we may choose
\(x_\alpha\in X^{\alpha+1}\setminus X^\alpha\) for each \(\alpha<\kappa\).
If \(\alpha<\beta<\kappa\), then \(x_\alpha\in X^{\alpha+1}\subseteq
X^\beta\) while \(x_\beta\notin X^\beta\), so \(x_\alpha\neq x_\beta\).
Hence \(\alpha\mapsto x_\alpha\) injects \(\kappa\) into \(B\), which is
impossible since \(\kappa>|B|\). So there is \(\alpha<\kappa\) with
\(X^\alpha=X^{\alpha+1}\); the least one is \(\operatorname{cl}(F)\), and
\(\operatorname{cl}(F)<\kappa\) means \(|\operatorname{cl}(F)|\leq|B|\).

\emph{Constancy.} Let \(X^\alpha=X^{\alpha+1}\). We show
\(X^\beta=X^\alpha\) for all \(\beta\geq\alpha\) by induction on
\(\beta\). For \(\beta=\alpha\) there is nothing to show. If
\(X^\beta=X^\alpha\), then \(X^{\beta+1}=F(X^\beta)=F(X^\alpha)=
X^{\alpha+1}=X^\alpha\). If \(\lambda>\alpha\) is a limit, then
\(X^\lambda=\bigcup_{\gamma<\lambda}X^\gamma\); the terms with
\(\gamma\geq\alpha\) equal \(X^\alpha\) by the induction hypothesis, and the
terms with \(\gamma<\alpha\) are contained in \(X^\alpha\) because \(F\) is
inductive. So \(X^\lambda=X^\alpha\). In particular
\(X^\infty=X^{\operatorname{cl}(F)}\) and
\(F(X^\infty)=X^{\operatorname{cl}(F)+1}=X^\infty\).

\emph{Monotone operators are inductive.} Let \(F\) be monotone. We prove by
induction on \(\beta\) the statement \(S(\beta)\): \(X^\alpha\subseteq
X^\beta\) for all \(\alpha\leq\beta\), and \(X^\beta\subseteq X^{\beta+1}\).
\begin{itemize}
 \item \(\beta=0\): \(X^0=\varnothing\) is contained in every set.
 \item \(\beta=\delta+1\): for \(\alpha\leq\delta\), \(S(\delta)\) gives
       \(X^\alpha\subseteq X^\delta\subseteq X^{\delta+1}\). Applying the
       monotone \(F\) to \(X^\delta\subseteq X^{\delta+1}\) gives
       \(X^{\delta+1}\subseteq X^{\delta+2}\).
 \item \(\beta\) a limit: for \(\alpha<\beta\), \(X^\alpha\subseteq
       X^\beta\) by the definition of \(X^\beta\) as a union. Also
       \(X^\alpha\subseteq X^{\alpha+1}=F(X^\alpha)\subseteq F(X^\beta)=
       X^{\beta+1}\), using \(S(\alpha)\) and monotonicity; taking the
       union over \(\alpha<\beta\) gives \(X^\beta\subseteq X^{\beta+1}\).
\end{itemize}
So the stages increase and \(F\) is inductive.
\end{proof}
\end{annexitem}

\begin{annexitem}{annex:p03-w03-p03-059-stages-lfp}{thm:p03-w03-p03-059-stages-lfp}{Stage construction of the extremal fixed points}
\begin{proof}
\emph{Item 1.} By \cref{thm:p03-w03-p03-059-closure-ordinal}, \(F\) is
inductive and \(X^\infty\) is a fixed point; hence
\(\operatorname{lfp}(F)\subseteq X^\infty\). For the converse let \(Z\)
satisfy \(F(Z)\subseteq Z\). By induction on \(\alpha\),
\(X^\alpha\subseteq Z\): \(X^0=\varnothing\subseteq Z\); if
\(X^\alpha\subseteq Z\) then \(X^{\alpha+1}=F(X^\alpha)\subseteq
F(Z)\subseteq Z\); and a union of subsets of \(Z\) is a subset of \(Z\).
Thus \(X^\infty\subseteq\bigcap\{Z:F(Z)\subseteq Z\}=\operatorname{lfp}(F)\)
by \cref{thm:p03-w03-p03-059-knaster-tarski}.

\emph{Item 3.} If \(X\subseteq Y\), then \(\overline Y\subseteq\overline
X\), so \(F(\overline Y)\subseteq F(\overline X)\) and
\(F^d(X)=\overline{F(\overline X)}\subseteq\overline{F(\overline Y)}=
F^d(Y)\); thus \(F^d\) is monotone. Moreover \(F^d(X)=X\) holds exactly
when \(F(\overline X)=\overline X\), so complementation is a bijection
between the fixed points of \(F^d\) and those of \(F\), and it reverses
inclusion. It therefore sends the least fixed point of \(F^d\) to the
greatest fixed point of \(F\) and the greatest fixed point of \(F^d\) to
the least fixed point of \(F\):
\(\operatorname{gfp}(F)=\overline{\operatorname{lfp}(F^d)}\) and
\(\operatorname{lfp}(F)=\overline{\operatorname{gfp}(F^d)}\).

\emph{Item 2.} Write \(X_d^\alpha\) for the stages of \(F^d\). We show
\(Y^\alpha=\overline{X_d^\alpha}\) for all \(\alpha\). At \(0\),
\(Y^0=B=\overline\varnothing\). At successors, if
\(Y^\alpha=\overline{X_d^\alpha}\) then \(Y^{\alpha+1}=F(\overline{X_d^\alpha})=
\overline{F^d(X_d^\alpha)}=\overline{X_d^{\alpha+1}}\), since
\(F(\overline X)=\overline{F^d(X)}\) by the definition of \(F^d\). At limits,
\(\bigcap_{\alpha<\lambda}\overline{X_d^\alpha}=
\overline{\bigcup_{\alpha<\lambda}X_d^\alpha}\). Since \(F^d\) is
monotone, its stages increase and become constant at
\(\operatorname{lfp}(F^d)\) by item~1 and
\cref{thm:p03-w03-p03-059-closure-ordinal}. Taking complements, the
\(Y^\alpha\) decrease and become constant at
\(\overline{\operatorname{lfp}(F^d)}=\operatorname{gfp}(F)\), by item~3.
\end{proof}
\end{annexitem}

\begin{annexitem}{annex:p03-w03-p03-059-finite-stages}{thm:p03-w03-p03-059-finite-stages}{Finite stages reach the fixed points in polynomially many steps}
\begin{proof}
Put \(B=A^k\), so \(|B|=n^k\). By
\cref{thm:p03-w03-p03-059-closure-ordinal} the finite stages increase. If
\(X^m\neq X^{m+1}\) for all \(m<n^k\), then each of these inclusions is
strict, so \(|X^m|\geq m\) for \(m\leq n^k\), and
\(|X^{n^k}|\geq n^k\) forces \(X^{n^k}=B\) and \(X^{n^k+1}=B\). In every case
there is \(m\leq n^k\) with \(X^m=X^{m+1}\). By the constancy part of
\cref{thm:p03-w03-p03-059-closure-ordinal}, all later stages, including
\(X^\omega=\bigcup_{j<\omega}X^j\) and all transfinite ones, equal \(X^m\);
so \(X^\infty=X^m\), which is \(\operatorname{lfp}(F)\) by
\cref{thm:p03-w03-p03-059-stages-lfp}. The downward stages are the
complements of the stages of the monotone operator \(F^d\) (proof of
\cref{thm:p03-w03-p03-059-stages-lfp}, item~2), so the same bound gives
\(\operatorname{gfp}(F)=Y^m\) for some \(m\leq n^k\).

For the complexity claim, a subset of \(A^k\) is represented by its
characteristic vector of length \(n^k\), and ``\(F\) is computable in
polynomial time'' means that some polynomial \(q\) bounds the time to compute
\(F(X)\) from \(X\) for all \(X\subseteq A^k\). The algorithm starts from
the all-zero vector and applies \(F\) until two consecutive vectors agree.
By the first part this happens after at most \(n^k+1\) applications, each
followed by a comparison of two vectors of length \(n^k\); the total time
is at most \((n^k+1)(q(n)+c\,n^k)\) for a constant \(c\), which is
polynomial in \(n\). The greatest fixed point is computed in the same way
from the all-one vector.
\end{proof}
\end{annexitem}

\begin{annexitem}{annex:p03-w03-p03-059-positive-monotone}{thm:p03-w03-p03-059-positive-monotone}{Positivity gives monotone operators}
\begin{proof}
Formulas are taken with the connectives \(\neg,\land,\lor\), the
quantifiers \(\exists,\forall\), and (for LFP formulas) the constructors
\(\mathbf{lfp}\) and \(\mathbf{gfp}\); implications and biconditionals are
first rewritten. Fix a structure \(\mathfrak A\). Say that a formula
\(\chi\) is \emph{upward in \(R\)} if, for every assignment \(s\) to its
other free variables (first-order and relational) and all
\(S\subseteq S'\subseteq A^k\), \((\mathfrak A,S)\entails\chi[s]\) implies
\((\mathfrak A,S')\entails\chi[s]\), and \emph{downward in \(R\)} if the
same holds with the conclusion and hypothesis on \(S\) and \(S'\)
exchanged. We prove by induction on \(\chi\): if every occurrence of \(R\)
in \(\chi\) is positive, \(\chi\) is upward in \(R\); if every occurrence is
negative, \(\chi\) is downward in \(R\).
\begin{itemize}
 \item Atoms. \(R\bar s\) has one positive occurrence, and
       \(\bar s^{\mathfrak A}\in S\subseteq S'\) gives upwardness; the
       downward hypothesis does not apply. An atom without \(R\) has a truth
       value independent of \(S\) and is both upward and downward.
 \item \(\neg\chi\). Its occurrences of \(R\) are positive exactly when those
       of \(\chi\) are negative. If \(\chi\) is downward and \(S\subseteq S'\)
       with \((\mathfrak A,S)\entails\neg\chi[s]\), then
       \((\mathfrak A,S')\entails\chi[s]\) would give
       \((\mathfrak A,S)\entails\chi[s]\); so \(\neg\chi\) is upward. The
       other case is symmetric.
 \item \(\chi_1\land\chi_2\), \(\chi_1\lor\chi_2\), \(\exists x\chi\),
       \(\forall x\chi\). Polarity of each occurrence is inherited from the
       immediate subformula, and conjunction, disjunction and the
       quantifiers (applied for each value of \(x\)) preserve upwardness and
       downwardness.
 \item \(\theta=[\mathbf{lfp}\,P\bar y\,.\,\chi](\bar t)\) with \(P\neq R\).
       For each interpretation \(S\) of \(R\), let \(G_S\) be the operator
       of \(\chi\) in \(P\) (with the other variables at their assigned
       values); it is monotone because \(P\) occurs positively in \(\chi\).
       Suppose the occurrences of \(R\) in \(\chi\) are positive, so by the
       induction hypothesis \(\chi\) is upward in \(R\) and \(G_S(Z)
       \subseteq G_{S'}(Z)\) for all \(Z\) when \(S\subseteq S'\). Then
       \(L'\defeq\operatorname{lfp}(G_{S'})\) satisfies
       \(G_S(L')\subseteq G_{S'}(L')=L'\), so
       \(\operatorname{lfp}(G_S)\subseteq L'\) by
       \cref{thm:p03-w03-p03-059-knaster-tarski}; hence \(\theta\) is
       upward. For \(\mathbf{gfp}\), \(M\defeq\operatorname{gfp}(G_S)\)
       satisfies \(M=G_S(M)\subseteq G_{S'}(M)\), so
       \(M\subseteq\operatorname{gfp}(G_{S'})\). The negative case is the
       same with \(S\) and \(S'\) exchanged.
\end{itemize}
Applying the claim to \(\varphi\) and to the tuples \(\bar a\) gives
monotonicity (respectively antitonicity) of \(F_\varphi\).
\end{proof}
\end{annexitem}

\begin{annexitem}{annex:p03-w03-p03-059-gfp-dual}{thm:p03-w03-p03-059-gfp-dual}{Greatest fixed points through negated least fixed points}
\begin{proof}
\emph{Polarity.} An occurrence of \(R\) in \(\varphi\) lies under an even
number of negations. In \(\varphi[R/\neg R]\) the atom at that place is
preceded by one more negation, and in \(\neg\varphi[R/\neg R]\) by one more
again, so the count is even and every occurrence is positive. Hence the
right-hand side is an LFP formula.

\emph{Substitution.} For every interpretation \(X\subseteq A^k\) of \(R\)
and every assignment \(s\),
\[
 (\mathfrak A,X)\entails\varphi[R/\neg R][s]
 \iff
 (\mathfrak A,\overline X)\entails\varphi[s].
\]
This is proved by induction on \(\varphi\). For an atom \(R\bar s\), the
left side is \(\bar s^{\mathfrak A}\notin X\), that is
\(\bar s^{\mathfrak A}\in\overline X\). Atoms without \(R\) do not depend on
the interpretation of \(R\), and the Boolean and quantifier steps are
immediate. For a fixed-point subformula with fixed-point variable
\(P\neq R\), the induction hypothesis applied to its body shows that, for
every interpretation of \(P\), the body after substitution under \(X\) and
the original body under \(\overline X\) define the same set; so the two
operators coincide and have the same least and greatest fixed points.

\emph{Operators.} Write \(\varphi'\defeq\neg\varphi[R/\neg R]\). By the
substitution claim,
\[
\begin{aligned}
 F_{\varphi'}(X)&=\{\bar a:(\mathfrak A,X)\not\entails
   \varphi[R/\neg R](\bar a)\}\\
 &=\{\bar a:(\mathfrak A,\overline X)\not\entails\varphi(\bar a)\}\\
 &=\overline{F_\varphi(\overline X)}=F_\varphi^d(X).
\end{aligned}
\]
By \cref{thm:p03-w03-p03-059-stages-lfp}, item~3,
\(\operatorname{gfp}(F_\varphi)=\overline{\operatorname{lfp}(F_\varphi^d)}
=\overline{\operatorname{lfp}(F_{\varphi'})}\). So
\(\bar t^{\mathfrak A}\in\operatorname{gfp}(F_\varphi)\) exactly when
\(\bar t^{\mathfrak A}\notin\operatorname{lfp}(F_{\varphi'})\), which is the
displayed equivalence.
\end{proof}
\end{annexitem}

\begin{annexitem}{annex:p03-w03-p03-059-pos-basic}{thm:p03-w03-p03-059-pos-basic}{Monotonicity and substitution for positive fixed-point formulas}
\begin{proof}
Throughout, use the renaming convention at the start of this annex:
no first-order or relation variable is both free and bound in a formula
under discussion, and bound relation symbols avoid those of any replacement
formula. In particular, a symbol \(R\) being substituted is not rebound
inside \(\theta\). Renaming bound variables does not change truth.

\emph{Item~1, together with the well-definedness of the semantics.} We
prove by induction on the formation of \(\theta\in\mathrm{Pos}[\sigma]\)
that the truth of \(\theta\) is defined in every \(\sigma\)-structure under
every assignment, and that \(\theta\) is monotone in every \(P\in\sigma\)
that occurs in no literal \(\neg P\bar t\) of \(\theta\). Let
\(\mathfrak A,\mathfrak A'\) differ only in that
\(P^{\mathfrak A}\subseteq P^{\mathfrak A'}\).
\begin{itemize}
 \item Literals. \(P\bar t\) true in \(\mathfrak A\) is true in
       \(\mathfrak A'\) because the interpretation grew. A literal not
       built from \(P\) has the same truth value in both structures. The
       literal \(\neg P\bar t\) does not occur by hypothesis.
 \item \(\land\) and \(\lor\) preserve the implication componentwise. For
       \(\exists x\) and \(\forall x\), the two structures have the same
       universe, so a witness, or every value, for \(x\) in \(\mathfrak A\)
       serves in \(\mathfrak A'\).
 \item \(\theta=[\mathbf{lfp}\,R_j:\mathcal S](\bar t)\) with bodies
       \(\varphi_1,\ldots,\varphi_m\). By the induction hypothesis each
       \(\varphi_i\) has defined truth values. Since no literal
       \(\neg R_l\bar t\) occurs in any body, the induction hypothesis for
       the symbol \(R_l\), applied to the structures
       \((\mathfrak A,X_1,\ldots,X_m)\) and
       \((\mathfrak A,X_1,\ldots,X_l',\ldots,X_m)\) with
       \(X_l\subseteq X_l'\), shows that enlarging one argument of
       \(F_{\mathcal S}\) enlarges every component of its value. Enlarging
       the arguments one at a time, \(F_{\mathcal S}\) is monotone on
       \(\mathcal P(B)\). By \cref{thm:p03-w03-p03-059-knaster-tarski} it
       has a least fixed point, so the truth of \(\theta\) is defined.

       For monotonicity in \(P\): no literal \(\neg P\bar t\) occurs in the
       bodies, so the induction hypothesis, applied to
       \((\mathfrak A,\bar X)\) and \((\mathfrak A',\bar X)\), gives
       \(F^{\mathfrak A}_{\mathcal S}(\bar X)\subseteq
       F^{\mathfrak A'}_{\mathcal S}(\bar X)\) for every \(\bar X\) and
       every value of the parameters. Let
       \(L'\defeq\operatorname{lfp}(F^{\mathfrak A'}_{\mathcal S})\). Then
       \(F^{\mathfrak A}_{\mathcal S}(L')\subseteq
       F^{\mathfrak A'}_{\mathcal S}(L')=L'\), so \(L'\) belongs to
       \(\{X:F^{\mathfrak A}_{\mathcal S}(X)\subseteq X\}\), whose
       intersection is \(\operatorname{lfp}(F^{\mathfrak A}_{\mathcal S})\)
       by \cref{thm:p03-w03-p03-059-knaster-tarski}. Hence
       \(\operatorname{lfp}(F^{\mathfrak A}_{\mathcal S})\subseteq L'\), and
       if \(\bar t\) denotes a tuple of the \(j\)-th component of the first
       fixed point it denotes one of the second.
\end{itemize}

\emph{Item~2, syntax.} Replacing literals by formulas of
\(\mathrm{Pos}[\sigma']\) and keeping \(\land,\lor,\exists,\forall\)
yields formulas of \(\mathrm{Pos}[\sigma']\), by induction on \(\theta\). At
a fixed-point atom with system \(\mathcal S\), the substituted bodies lie in
\(\mathrm{Pos}[\sigma'\cup\{R_1,\ldots,R_m\}]\) by the induction
hypothesis, and they contain no literal \(\neg R_l\bar t\): their own
literals other than \(R\)-literals are unchanged, and \(\alpha,\beta\) do
not contain the bound symbols \(R_l\). So the substituted atom is formed by
item~3 of \cref{def:p03-w03-p03-059-simultaneous-lfp}.

\emph{Item~2, semantics.} Let \(\bar v\) be the free variables of
\(\alpha\) and \(\beta\) other than \(\bar u\). After renaming, no variable
of \(\bar v\) is bound in \(\theta\). We prove, by induction on \(\theta\),
that for every assignment \(s'\) agreeing with \(s\) on \(\bar v\),
\(\mathfrak B\entails\theta[R:\alpha/\beta][s']\) if and only if
\(\mathfrak B[R\mapsto Q]\entails\theta[s']\). Since \(\alpha\) and
\(\beta\) depend only on the values of \(\bar u\) and \(\bar v\), \(Q\) and
the hypothesis on \(\beta\) are the same under \(s'\) as under \(s\).
\begin{itemize}
 \item \(\theta=R\bar t\). Then \(\mathfrak B\entails\alpha(\bar t)[s']\)
       if and only if \(s'(\bar t)\in Q\), if and only if
       \(\mathfrak B[R\mapsto Q]\entails R\bar t[s']\).
 \item \(\theta=\neg R\bar t\). Then \(\mathfrak B\entails\beta(\bar t)[s']\)
       if and only if \(s'(\bar t)\notin Q\), by the hypothesis on
       \(\beta\). This is the only place where that hypothesis is used, so
       it may be dropped when no such literal occurs.
 \item Other literals do not contain \(R\), and \(\mathfrak B\) and
       \(\mathfrak B[R\mapsto Q]\) agree on their symbols.
 \item \(\land,\lor,\exists x,\forall x\): immediate from the induction
       hypothesis, since \(x\notin\bar v\) and so every modified
       assignment still agrees with \(s\) on \(\bar v\).
 \item \(\theta=[\mathbf{lfp}\,R_j:\mathcal S](\bar t)\). The symbols
       \(R_1,\ldots,R_m\) of \(\mathcal S\) differ from \(R\) and do not
       occur in \(\alpha,\beta\). Fix interpretations \(\bar X\) of them and
       values of the parameters. In \((\mathfrak B,\bar X)\) the formulas
       \(\alpha\) and \(\beta\) define the same sets as in \(\mathfrak B\),
       so the induction hypothesis for the body \(\varphi_i\), read over
       the vocabulary \(\sigma\cup\{R_1,\ldots,R_m\}\), gives
       \((\mathfrak B,\bar X)\entails\varphi_i[R:\alpha/\beta](\bar a)\) if
       and only if \((\mathfrak B[R\mapsto Q],\bar X)\entails
       \varphi_i(\bar a)\). Thus the substituted system has in
       \(\mathfrak B\) the same operator as \(\mathcal S\) has in
       \(\mathfrak B[R\mapsto Q]\), hence the same least fixed point, and
       the two atoms have the same truth value.\qedhere
\end{itemize}
\end{proof}
\end{annexitem}

\begin{annexitem}{annex:p03-w03-p03-059-stage-comparison}{thm:p03-w03-p03-059-stage-comparison}{Positive stage comparison and positive negation of least fixed points}
\begin{proof}
\emph{Facts about the stages.} Fix a finite \(\sigma\)-structure
\(\mathfrak A\) and values of \(\bar p\), and write \(|\bar a|\) for
\(|\bar a|_F\). By \cref{thm:p03-w03-p03-059-closure-ordinal} and
\cref{thm:p03-w03-p03-059-finite-stages}, the stages increase, there is a
least natural number \(c\) with \(X^{c+1}=X^c\), and
\(X^m=X^c=\operatorname{lfp}(F)\) for every \(m\geq c\). For \(m<c\) the
inclusion \(X^m\subseteq X^{m+1}\) is proper, by the choice of \(c\).
\begin{enumerate}
 \item[(F1)] \(X^m=\{\bar a:|\bar a|\leq m\}\) for every \(m\). If
       \(\bar a\in X^m\), then \(\bar a\in X^c=\operatorname{lfp}(F)\), so
       \(|\bar a|\) is finite and at most \(m\). If \(|\bar a|=r\leq m\),
       then \(\bar a\in X^r\subseteq X^m\).
 \item[(F2)] Every finite rank lies in \(\{1,\ldots,c\}\), since
       \(X^0=\varnothing\) and \(\operatorname{lfp}(F)=X^c\); and every
       \(r\in\{1,\ldots,c\}\) is the rank of some tuple, namely of any
       element of \(X^r\setminus X^{r-1}\neq\varnothing\).
 \item[(F3)] If \(|\bar c|=i<\infty\), then
       \(\{\bar a:\bar a\le_F\bar c\}=X^i\) and
       \(\{\bar a:\bar c<_F\bar a\}=A^k\setminus X^i\), by (F1).
 \item[(F4)] For every \(X\subseteq A^k\),
       \((\mathfrak A,X)\entails\theta^+(\bar a)\) if and only if
       \(\bar a\in F(X)\), and \((\mathfrak A,X)\entails\theta^-(\bar a)\)
       if and only if \(\bar a\notin F(X)\). In particular, by
       \cref{thm:p03-w03-p03-059-pos-basic}, item~2, with \(\alpha=\bot\)
       and \(\beta=\top\), the formula \(\theta^+(\bar a)[R:\varnothing]\)
       holds if and only if \(\bar a\in F(\varnothing)=X^1\), and
       \(\theta^-(\bar a)[R:\varnothing]\) if and only if
       \(\bar a\notin X^1\).
\end{enumerate}

\emph{The system.} Let \(L,M\) be new relation symbols of arity \(2k\),
let \(\bar y,\bar z\) be fresh \(k\)-tuples of variables, and for a
\(k\)-tuple \(\bar z\) put
\[
 \theta^+(\bar x)\langle\bar z\rangle\defeq
 \theta^+[R:L(\cdot,\bar z)/M(\bar z,\cdot)],\qquad
 \theta^-(\bar y)\langle\bar z\rangle\defeq
 \theta^-(\bar y)[R:L(\cdot,\bar z)/M(\bar z,\cdot)],
\]
where \(L(\cdot,\bar z)\) is \(L\bar u\bar z\), \(M(\bar z,\cdot)\) is
\(M\bar z\bar u\), and \(\theta^-(\bar y)\) is \(\theta^-\) with \(\bar x\)
renamed to \(\bar y\). Let \(\mathcal S_{\mathrm{cmp}}\) be the system
\[
\begin{aligned}
 L\bar x\bar y&\defeq\theta^+(\bar x)[R:\varnothing]\lor
   \exists\bar z\bigl(M\bar z\bar y\land\theta^+(\bar x)\langle\bar z\rangle\bigr),\\
 M\bar x\bar y&\defeq\bigl(\theta^+(\bar x)[R:\varnothing]\land
   \theta^-(\bar y)[R:\varnothing]\bigr)\lor
   \exists\bar z\bigl(L\bar z\bar z\land\theta^+(\bar x)\langle\bar z\rangle
   \land\theta^-(\bar y)\langle\bar z\rangle\bigr).
\end{aligned}
\]
By \cref{thm:p03-w03-p03-059-pos-basic}, item~2, both bodies lie in
\(\mathrm{Pos}[\sigma\cup\{L,M\}]\). The symbols \(L\) and \(M\) occur in
them only in the unnegated atoms written out and in the unnegated atoms
\(L\bar u\bar z\), \(M\bar z\bar u\) put in for \(R\)-literals, so no
literal \(\neg L\bar t\) or \(\neg M\bar t\) occurs. The parameters are
\(\bar p\). Put
\[
 \mathrm{le}(\bar u,\bar w)\defeq[\mathbf{lfp}\,L:\mathcal S_{\mathrm{cmp}}](\bar u,\bar w),
 \qquad
 \mathrm{lt}(\bar u,\bar w)\defeq[\mathbf{lfp}\,M:\mathcal S_{\mathrm{cmp}}](\bar u,\bar w),
\]
which belong to \(\mathrm{Pos}[\sigma]\). This system is a variant of the
one displayed by Immerman \cite[p.~98]{Immerman1986}: the complement of a
stage is named by \(M\bar z\bar u\), and the second disjunct of \(M\)
quantifies over tuples of finite rank.

\emph{The stages of the system.} Let \((L^j,M^j)\) be the stages of
\(F_{\mathcal S_{\mathrm{cmp}}}\), which is monotone by
\cref{thm:p03-w03-p03-059-pos-basic}, item~1. We show by induction on
\(j\) that
\[
 L^j=\{(\bar a,\bar b):|\bar a|\leq j,\ |\bar a|\leq|\bar b|\},\qquad
 M^j=\{(\bar a,\bar b):|\bar a|\leq j,\ |\bar a|<|\bar b|\}.
\]
For \(j=0\) both sides are empty, since finite ranks are at least \(1\).
Assume the claim for \(j\) and let \(\mathfrak A_j\defeq(\mathfrak
A,L^j,M^j)\).

First, let \(|\bar c|=i\leq j\). In \(\mathfrak A_j\) the formula
\(L\bar u\bar c\) defines \(\{\bar a:|\bar a|\leq j,\ |\bar a|\leq i\}=X^i\)
by (F1), and \(M\bar c\bar u\) defines
\(\{\bar a:i<|\bar a|\}=A^k\setminus X^i\). By
\cref{thm:p03-w03-p03-059-pos-basic}, item~2, with \(Q=X^i\), and by
(F4), since \(L,M\) do not occur in \(\theta^\pm\),
\[
 \mathfrak A_j\entails\theta^+(\bar a)\langle\bar c\rangle\iff
 \bar a\in F(X^i)=X^{i+1},\qquad
 \mathfrak A_j\entails\theta^-(\bar b)\langle\bar c\rangle\iff
 \bar b\notin X^{i+1}.
\]
In the body of \(L\) the quantified \(\bar z\) must satisfy
\(M\bar z\bar y\), and in the body of \(M\) it must satisfy
\(L\bar z\bar z\); by the induction hypothesis either condition gives
\(|\bar z|\leq j\), so the displayed equivalences apply to every witness.
Together with (F4) this yields
\[
\begin{aligned}
 L^{j+1}&=\{(\bar a,\bar b):|\bar a|=1\}\cup
   \{(\bar a,\bar b):\exists\bar c\,(|\bar c|\leq j,\ |\bar c|<|\bar b|,\
   |\bar a|\leq|\bar c|+1)\},\\
 M^{j+1}&=\{(\bar a,\bar b):|\bar a|=1<|\bar b|\}\cup
   \{(\bar a,\bar b):\exists\bar c\,(|\bar c|\leq j,\
   |\bar a|\leq|\bar c|+1<|\bar b|)\},
\end{aligned}
\]
where \(\bar a\in X^1\) is \(|\bar a|=1\), and \(\bar a\in X^{i+1}\) is
\(|\bar a|\leq i+1\), by (F1).

The first set equals \(\{(\bar a,\bar b):|\bar a|\leq j+1,\
|\bar a|\leq|\bar b|\}\). If \(|\bar a|=1\), then \(|\bar a|\leq j+1\) and
\(|\bar a|\leq|\bar b|\) because every rank is at least \(1\). If
\(\bar c\) is a witness of the second part, then
\(|\bar a|\leq|\bar c|+1\leq j+1\), and \(|\bar c|<|\bar b|\) gives
\(|\bar c|+1\leq|\bar b|\), ranks being natural numbers or \(\infty\).
Conversely let \(|\bar a|=r\leq j+1\) and \(r\leq|\bar b|\). If \(r=1\),
the first part applies. If \(r\geq2\), then \(r\leq c\) by (F2), so by (F2)
some \(\bar c\) has rank \(r-1\); it satisfies \(|\bar c|=r-1\leq j\),
\(|\bar c|<r\leq|\bar b|\) and \(|\bar a|=|\bar c|+1\).

The second set equals \(\{(\bar a,\bar b):|\bar a|\leq j+1,\
|\bar a|<|\bar b|\}\). If \(|\bar a|=1<|\bar b|\) this is clear. A witness
\(\bar c\) of the second part gives \(|\bar a|\leq|\bar c|+1\leq j+1\) and
\(|\bar a|\leq|\bar c|+1<|\bar b|\). Conversely, if \(|\bar a|=r\leq j+1\)
and \(r<|\bar b|\), then for \(r=1\) the first part applies, and for
\(r\geq2\) a tuple \(\bar c\) of rank \(r-1\) gives \(|\bar c|\leq j\) and
\(|\bar a|=|\bar c|+1<|\bar b|\). This proves the claim.

\emph{The fixed point.} By \cref{thm:p03-w03-p03-059-closure-ordinal} and
\cref{thm:p03-w03-p03-059-stages-lfp}, item~1, applied on \(\mathcal P(B)\)
for the system, the stages increase and their union is
\(\operatorname{lfp}(F_{\mathcal S_{\mathrm{cmp}}})\); on the finite set
\(B\) this union is reached at a finite stage. The union of the sets in the
claim is \(\{(\bar a,\bar b):|\bar a|<\infty,\ |\bar a|\leq|\bar b|\}\),
which is \(\le_F\), and \(\{(\bar a,\bar b):|\bar a|<|\bar b|\}\), which
is \(<_F\). So \(\mathrm{le}\) and \(\mathrm{lt}\) define \(\le_F\) and
\(<_F\).

\emph{Membership.} Put \(\mathrm{in}(\bar u)\defeq\mathrm{le}(\bar u,\bar
u)\). It holds exactly when \(|\bar u|<\infty\), that is, when \(\bar u\in
\operatorname{lfp}(F)\).

\emph{Maximal rank.} Put
\[
 \mathrm{MAX}(\bar z)\defeq\mathrm{le}(\bar z,\bar z)\land
 \forall\bar y\bigl(\mathrm{le}(\bar y,\bar z)\lor
 \theta^-(\bar y)[R:\mathrm{le}(\cdot,\bar z)/\mathrm{lt}(\bar z,\cdot)]\bigr),
\]
a formula of \(\mathrm{Pos}[\sigma]\) by
\cref{thm:p03-w03-p03-059-pos-basic}, item~2. We show that
\(\mathrm{MAX}(\bar c)\) holds exactly when \(|\bar c|=c\). The first
conjunct says that \(i\defeq|\bar c|\) is finite. Then, by (F3),
\(\mathrm{le}(\cdot,\bar c)\) defines \(X^i\) and
\(\mathrm{lt}(\bar c,\cdot)\) its complement, so by
\cref{thm:p03-w03-p03-059-pos-basic}, item~2, and (F4) the substituted
\(\theta^-(\bar b)\) holds exactly when \(\bar b\notin F(X^i)=X^{i+1}\).
Hence the second conjunct says that every \(\bar b\) lies in \(X^i\) or
outside \(X^{i+1}\), that is, \(X^{i+1}\subseteq X^i\). As the stages
increase, this means \(X^{i+1}=X^i\), which holds exactly when \(i\geq c\).
Since \(i\leq c\) by (F2), it holds exactly when \(i=c\).

\emph{The complement.} Put
\[
 \mathrm{out}(\bar u)\defeq\forall\bar y\,\theta^-(\bar y)[R:\varnothing]
 \lor\exists\bar z\bigl(\mathrm{MAX}(\bar z)\land\mathrm{lt}(\bar z,\bar u)\bigr).
\]
If \(\operatorname{lfp}(F)=\varnothing\), then
\(X^1=F(\varnothing)\subseteq F(\operatorname{lfp}(F))=\varnothing\), so the
first disjunct holds by (F4), and \(\bar u\notin\operatorname{lfp}(F)\)
for every \(\bar u\). If \(\operatorname{lfp}(F)\neq\varnothing\), then
\(X^1\neq\varnothing\), since otherwise \(\varnothing\) would be a fixed
point and so equal to the least one; so the first disjunct fails. Then
\(c\geq1\), a tuple of rank \(c\) exists by (F2), and the tuples
satisfying \(\mathrm{MAX}\) are exactly those of rank \(c\). For such a
tuple \(\bar z\), \(\mathrm{lt}(\bar z,\bar u)\) says \(c<|\bar u|\), that
is, \(\bar u\notin X^c=\operatorname{lfp}(F)\) by (F1). In both cases
\(\mathrm{out}(\bar u)\) holds exactly when
\(\bar u\notin\operatorname{lfp}(F)\).

All four formulas lie in \(\mathrm{Pos}[\sigma]\), and their free variables
are among \(\bar u,\bar w,\bar p\).
\end{proof}
\end{annexitem}

\begin{annexitem}{annex:p03-w03-p03-059-immerman-normal-form}{thm:p03-w03-p03-059-immerman-normal-form}{Immerman's normal form on finite structures}
\begin{proof}
All equivalences below are on finite structures, under every assignment
and every interpretation of the relation symbols of the vocabulary in
question. Implication and the biconditional are read as abbreviations in
terms of \(\neg,\land,\lor\), as in the positivity condition of
\cref{def:p03-w03-p03-059-lfp}.

\emph{Item~1.} By induction on the LFP formula \(\chi\), simultaneously for
all relational vocabularies \(\sigma\).
\begin{itemize}
 \item Atomic \(\chi\) (a relation atom or an equation): \(\chi^+\defeq\chi\)
       and \(\chi^-\defeq\neg\chi\), both literals.
 \item \((\neg\chi_1)^+\defeq\chi_1^-\) and
       \((\neg\chi_1)^-\defeq\chi_1^+\).
 \item \((\chi_1\land\chi_2)^+\defeq\chi_1^+\land\chi_2^+\) and
       \((\chi_1\land\chi_2)^-\defeq\chi_1^-\lor\chi_2^-\); dually for
       \(\lor\). \((\exists x\,\chi_1)^+\defeq\exists x\,\chi_1^+\) and
       \((\exists x\,\chi_1)^-\defeq\forall x\,\chi_1^-\); dually for
       \(\forall\). The equivalences follow from those for \(\chi_1,\chi_2\)
       and De Morgan's laws.
 \item \(\chi=[\mathbf{lfp}\,R\bar x\,.\,\varphi](\bar t)\), where
       \(\varphi\) is an LFP formula of vocabulary \(\sigma\cup\{R\}\) with
       \(R\) positive, and \(\bar p\) are its parameters. The induction
       hypothesis, applied over \(\sigma\cup\{R\}\), gives
       \(\varphi^+,\varphi^-\in\mathrm{Pos}[\sigma\cup\{R\}]\) with free
       variables among \(\bar x,\bar p\), equivalent to \(\varphi\) and
       \(\neg\varphi\). On a finite \(\sigma\)-structure the operator
       \(X\mapsto\{\bar a:(\mathfrak A,X)\entails\varphi^+(\bar a)\}\) is
       \(F_\varphi\), which is monotone by
       \cref{thm:p03-w03-p03-059-positive-monotone}. So
       \cref{thm:p03-w03-p03-059-stage-comparison} applies to
       \(\theta^\pm\defeq\varphi^\pm\), and we put
       \(\chi^+\defeq\mathrm{in}(\bar t)\) and
       \(\chi^-\defeq\mathrm{out}(\bar t)\), with \(\bar t\) substituted
       for \(\bar u\). By the semantics of \cref{def:p03-w03-p03-059-lfp},
       \(\chi\) holds exactly when \(\bar t\) denotes a tuple of
       \(\operatorname{lfp}(F_\varphi)\). The free variables are among
       \(\bar t,\bar p\), the free variables of \(\chi\).
 \item \(\chi=[\mathbf{gfp}\,R\bar x\,.\,\varphi](\bar t)\), with
       \(\varphi^\pm\) as before. Put
       \(\theta^+\defeq(\varphi^-)^{\sim}\) and
       \(\theta^-\defeq(\varphi^+)^{\sim}\). In a structure
       \((\mathfrak A,X)\) the formula \(\neg R\bar u\) defines
       \(A^k\setminus X\) and \(R\bar u\) its complement, so
       \cref{thm:p03-w03-p03-059-pos-basic}, item~2, with
       \(\sigma'=\sigma\cup\{R\}\), gives
       \[
        (\mathfrak A,X)\entails\theta^+(\bar a)\iff
        (\mathfrak A,A^k\setminus X)\entails\varphi^-(\bar a)\iff
        \bar a\notin F_\varphi(A^k\setminus X)\iff
        \bar a\in F^d_\varphi(X),
       \]
       and in the same way \((\mathfrak A,X)\entails\theta^-(\bar a)\) if
       and only if \(\bar a\notin F^d_\varphi(X)\). So \(\theta^-\) is
       equivalent to \(\neg\theta^+\), and the operator of \(\theta^+\) is
       \(F^d_\varphi\), which is monotone with
       \(\operatorname{lfp}(F^d_\varphi)=A^k\setminus
       \operatorname{gfp}(F_\varphi)\) by
       \cref{thm:p03-w03-p03-059-stages-lfp}, item~3. Applying
       \cref{thm:p03-w03-p03-059-stage-comparison}, put
       \(\chi^+\defeq\mathrm{out}(\bar t)\) and
       \(\chi^-\defeq\mathrm{in}(\bar t)\).
\end{itemize}
The remaining step is to turn formulas of \(\mathrm{Pos}[\sigma]\) into
formulas of \(\mathrm{LFP}_1\). Call a system \emph{parameter-free} if the
free variables of each body \(\varphi_i\) are among \(\bar x_i\), and
\emph{first-order} if no body contains a fixed-point atom.

\emph{Item~2, step A: parameters.} We show that every
\(\theta\in\mathrm{Pos}[\sigma]\) is equivalent to a formula of
\(\mathrm{Pos}[\sigma]\) with the same free variables all of whose systems,
at every depth, are parameter-free. The proof is by induction on the
nesting depth of fixed-point atoms; literals are kept, and \(\land,\lor,
\exists,\forall\) are handled componentwise. Let
\(\theta=[\mathbf{lfp}\,R_j:\mathcal S](\bar t)\) and let \(\bar p\) be the
\(q\) parameters of \(\mathcal S\). Let \(\mathcal S^{\bar p}\) have new
symbols \(R'_i\) of arity \(k_i+q\), variable tuples \(\bar x_i\bar p\),
and bodies \(\varphi'_i\) obtained from \(\varphi_i\) by replacing every
literal \(R_l\bar s\) by \(R'_l\bar s\bar p\) (there are no literals
\(\neg R_l\bar s\)). The variables \(\bar p\) are free in \(\varphi_i\),
hence not bound there, so the free variables of \(\varphi'_i\) are among
\(\bar x_i\bar p\). Put
\(\theta'\defeq[\mathbf{lfp}\,R'_j:\mathcal S^{\bar p}](\bar t\bar p)\),
which has the same free variables as \(\theta\), and apply the induction
hypothesis to the bodies \(\varphi'_i\), whose depth is smaller; this
replaces them by equivalent bodies with the same free variables, so the
operator of \(\mathcal S^{\bar p}\) is unchanged and \(\mathcal
S^{\bar p}\) stays parameter-free.

To see \(\theta'\equiv\theta\), fix a finite \(\mathfrak A\) and values
\(\bar b\) of \(\bar p\). For \(X'\subseteq A^{k+q}\) write
\(X'|_{\bar b}\defeq\{\bar a:\bar a\bar b\in X'\}\). Applying
\cref{thm:p03-w03-p03-059-pos-basic}, item~2, to \(R_1,\ldots,R_m\) in turn,
with \(\alpha_l(\bar u)\defeq R'_l\bar u\bar p\) defining
\(X'_l|_{\bar b}\),
\[
 (\mathfrak A,\bar X')\entails\varphi'_i(\bar a\bar b)\iff
 (\mathfrak A,X'_1|_{\bar b},\ldots,X'_m|_{\bar b})\entails
 \varphi_i(\bar a)\ \text{with }\bar p\mapsto\bar b .
\]
So the \(\bar b\)-slice of \(F_{\mathcal S^{\bar p}}(\bar X')\) is
\(F_{\mathcal S}\), at parameters \(\bar b\), applied to the
\(\bar b\)-slices of \(\bar X'\). By induction on \(r\), the \(\bar b\)-slice
of the \(r\)-th stage of \(F_{\mathcal S^{\bar p}}\) is the \(r\)-th stage
of \(F_{\mathcal S}\) at \(\bar b\). Both operators are monotone, and by
\cref{thm:p03-w03-p03-059-closure-ordinal} on the finite sets involved both
reach their least fixed points at a common finite stage, by
\cref{thm:p03-w03-p03-059-stages-lfp}, item~1. Hence the \(\bar b\)-slice
of \(\operatorname{lfp}(F_{\mathcal S^{\bar p}})\) is
\(\operatorname{lfp}(F_{\mathcal S})\) at \(\bar b\), and \(\theta'\) holds
exactly when \(\theta\) does.

\emph{Item~2, step B: nested fixed points (transitivity).} We use
Beki\'c's principle. Let \(G:\mathcal P(U)\times\mathcal P(V)\to\mathcal
P(U)\) and \(H:\mathcal P(U)\times\mathcal P(V)\to\mathcal P(V)\) be
monotone, and for \(X\subseteq U\) let \(\mu(X)\) be the least fixed point
of \(Y\mapsto H(X,Y)\). Then \(\mu\) is monotone, and if \(X^*\) is the
least fixed point of \(X\mapsto G(X,\mu(X))\), the least fixed point of
\((X,Y)\mapsto(G(X,Y),H(X,Y))\) on \(\mathcal P(U)\times\mathcal P(V)\)
(identified with \(\mathcal P(U\sqcup V)\)) is \((X^*,\mu(X^*))\).

\emph{Proof of the principle.} If \(X\subseteq X'\), then
\(H(X,\mu(X'))\subseteq H(X',\mu(X'))=\mu(X')\), so \(\mu(X')\) is closed
under \(H(X,\cdot)\) and contains \(\mu(X)\) by
\cref{thm:p03-w03-p03-059-knaster-tarski}. Hence \(X\mapsto G(X,\mu(X))\) is
monotone and \(X^*\) exists. The pair \((X^*,\mu(X^*))\) is a fixed point
of \((G,H)\), so it contains the least fixed point \((X_0,Y_0)\).
Conversely, \(Y_0=H(X_0,Y_0)\) is a fixed point of \(H(X_0,\cdot)\), so
\(\mu(X_0)\subseteq Y_0\); then
\(G(X_0,\mu(X_0))\subseteq G(X_0,Y_0)=X_0\), so \(X^*\subseteq X_0\) by
\cref{thm:p03-w03-p03-059-knaster-tarski}, and
\(\mu(X^*)\subseteq\mu(X_0)\subseteq Y_0\). So the two pairs are equal.

We show that a formula \(\theta\in\mathrm{Pos}[\sigma]\) whose systems are
all parameter-free is equivalent to a formula built by \(\land,\lor,
\exists,\forall\) from literals and atoms of parameter-free first-order
systems. Let \(N(\theta)\) be the number of occurrences of fixed-point atoms
in \(\theta\) that lie inside the body of another fixed-point atom; we use
induction on \(N(\theta)\). If \(N(\theta)=0\), every atom is top-level and
its bodies contain no atoms, so \(\theta\) already has the required form.
If \(N(\theta)>0\), choose a top-level occurrence of an atom
\(\alpha=[\mathbf{lfp}\,R_j:\mathcal S](\bar t)\) one of whose bodies
\(\varphi_i\) contains an atom, and in \(\varphi_i\) an occurrence of an atom
\(\beta=[\mathbf{lfp}\,T_l:\mathcal T](\bar s)\) not inside another atom
of \(\varphi_i\). The system \(\mathcal T\), over
\(\sigma\cup\{R_1,\ldots,R_m\}\), has symbols \(T_1,\ldots,T_{m'}\), which
we may take distinct from all other symbols, and bodies \(\psi_h\). Let
\(\mathcal S^+\) be the system with symbols \(R_1,\ldots,R_m,
T_1,\ldots,T_{m'}\) whose bodies are those of \(\mathcal S\), except that in
\(\varphi_i\) the chosen occurrence of \(\beta\) is replaced by
\(T_l\bar s\), together with the bodies \(\psi_h\) of \(\mathcal T\). It is
a legitimate system: no literal \(\neg R_l\bar t\) occurs, since none
occurred in \(\varphi_i\), which contained the \(\psi_h\); no literal
\(\neg T_h\bar t\) occurs, by the formation of \(\beta\) and since \(T_l\bar
s\) is unnegated; and it is parameter-free, because \(\beta\), being
parameter-free, has exactly the free variables \(\bar s\), as \(T_l\bar s\)
does.

Fix \(\mathfrak A\). Let \(G\) be the \(R\)-part and \(H\) the \(T\)-part of
\(F_{\mathcal S^+}\), both monotone by
\cref{thm:p03-w03-p03-059-pos-basic}, item~1. For each \(\bar X\),
\(\mu(\bar X)\) is the least fixed point of \(\mathcal T\) in
\((\mathfrak A,\bar X)\), so in \((\mathfrak A,\bar X,\mu(\bar X))\) the
formulas \(\beta\) and \(T_l\bar s\) are equivalent under every assignment.
Since the chosen occurrence lies in \(\varphi_i\) only under connectives and
quantifiers, replacing it by an equivalent formula keeps the truth value,
so \(G(\bar X,\mu(\bar X))=F_{\mathcal S}(\bar X)\). By Beki\'c's principle
the \(R\)-part of \(\operatorname{lfp}(F_{\mathcal S^+})\) is
\(\operatorname{lfp}(F_{\mathcal S})\), so \(\alpha\) is equivalent to
\(\alpha^+\defeq[\mathbf{lfp}\,R_j:\mathcal S^+](\bar t)\). Replacing the
top-level occurrence \(\alpha\) by \(\alpha^+\) in \(\theta\) gives an
equivalent formula whose systems are parameter-free and in which \(N\) is
smaller by one: the occurrence \(\beta\) was inside a body and is gone,
the atoms inside the \(\psi_h\) are still inside a body, and no other
occurrence changes its position. This is the form taken by Immerman's
transitivity lemma \cite[Lemma~4.2, p.~96]{Immerman1986} when the inner
fixed point may depend on the outer one.

\emph{Item~2, step C1: simultaneous inductions on structures with at least
two elements.} Let \(\mathcal S\) be parameter-free and first-order with
\(m\geq2\) components. Its bodies are first-order formulas in which
negation stands only before atoms and no \(R_l\) is negated. Let \(U\) be
a new symbol of arity \(2+k_1+k_2\), let \(v_1,v_2,\bar y_1,\bar y_2\) be
fresh variables with \(|\bar y_i|=k_i\), and put
\[
 D_1(\bar u)\defeq\exists v\,\exists\bar y_2\,Uvv\bar u\bar y_2,\qquad
 D_2(\bar u)\defeq\exists v_1\exists v_2\exists\bar y_1\bigl(\neg v_1=v_2\land
 Uv_1v_2\bar y_1\bar u\bigr).
\]
Let \(\varphi_i^\natural\) be \(\varphi_i\) with \(R_1\) and \(R_2\)
replaced by \(D_1\) and \(D_2\). Let \(\mathcal S'\) have the symbols
\(U,R_3,\ldots,R_m\) with
\[
 Uv_1v_2\bar y_1\bar y_2\defeq\bigl(v_1=v_2\land\varphi_1^\natural(\bar y_1)\bigr)
 \lor\bigl(\neg v_1=v_2\land\varphi_2^\natural(\bar y_2)\bigr),
 \qquad R_i\bar x_i\defeq\varphi_i^\natural\ (i\geq3),
\]
where \(\varphi_i(\bar y_i)\) is \(\varphi_i\) with \(\bar x_i\) renamed. It
is parameter-free, first-order, and has no negated \(U\) or \(R_i\). Let
\(\mathfrak A\) have at least two elements and let \((R_1^r,\ldots,R_m^r)\)
and \((U^r,R_3'^r,\ldots,R_m'^r)\) be the stages of the two systems. We
show by induction on \(r\) that \(R_i'^r=R_i^r\) for \(i\geq3\) and
\[
 U^r=\{(a,a,\bar b,\bar c):\bar b\in R_1^r\}\cup
 \{(a,a',\bar b,\bar c):a\neq a',\ \bar c\in R_2^r\}.
\]
For \(r=0\) all sets are empty. Assume the claim for \(r\). In
\((\mathfrak A,U^r,R_3^r,\ldots,R_m^r)\) the formula \(D_1\) defines
\(R_1^r\), because \(A\) and \(A^{k_2}\) are nonempty, and \(D_2\) defines
\(R_2^r\), because \(A\) has two distinct elements. By
\cref{thm:p03-w03-p03-059-pos-basic}, item~2 (no negated \(R_1,R_2\)
occur, so no \(\beta\) is needed), \(\varphi_i^\natural\) holds there at
\(\bar a\) exactly when \(\varphi_i\) holds at \(\bar a\) in
\((\mathfrak A,R_1^r,\ldots,R_m^r)\), that is, when
\(\bar a\in R_i^{r+1}\). Reading off the bodies gives the claim for
\(r+1\). The stages of both systems become constant from some finite stage
on (\cref{thm:p03-w03-p03-059-closure-ordinal}), where they are the least
fixed points (\cref{thm:p03-w03-p03-059-stages-lfp}, item~1). Hence, on
structures with at least two elements, \([\mathbf{lfp}\,R_1:\mathcal
S](\bar t)\) is equivalent to \(\exists v\,\exists\bar y_2\,
[\mathbf{lfp}\,U:\mathcal S'](v,v,\bar t,\bar y_2)\),
\([\mathbf{lfp}\,R_2:\mathcal S](\bar t)\) to
\(\exists v_1\exists v_2\exists\bar y_1(\neg v_1=v_2\land
[\mathbf{lfp}\,U:\mathcal S'](v_1,v_2,\bar y_1,\bar t))\), and
\([\mathbf{lfp}\,R_i:\mathcal S](\bar t)\) to
\([\mathbf{lfp}\,R_i:\mathcal S'](\bar t)\) for \(i\geq3\). This is
Immerman's simultaneous induction lemma in his variant without constants,
where a tag is a pair of equal or unequal elements
\cite[Lemma~4.1, pp.~95--96]{Immerman1986}.

Repeating this \(m-1\) times, each atom of a parameter-free first-order
system is equivalent, on structures with at least two elements, to a
formula built with \(\exists\), \(\land\) and equality literals from atoms
\([\mathbf{lfp}\,U:\mathcal S_1](\bar t')\) of a single-component system
\(U\bar z\defeq\varphi\), with \(\varphi\) first-order, positive in \(U\)
and with free variables among \(\bar z\). Such an atom is equivalent to
\(\exists\bar z\bigl(\bigwedge_i z_i=t'_i\land
[\mathbf{lfp}\,U\bar z\,.\,\varphi](\bar z)\bigr)\), a formula of
\(\mathrm{LFP}_1\) in the sense of \cref{def:p03-w03-p03-059-lfp1}.

\emph{Item~2, step C2: one-element structures.} Let \(\mathcal S\) be
parameter-free and first-order with \(m\) components. Put
\(\varphi_i^{(0)}\defeq\bot\) and let \(\varphi_i^{(r+1)}\) be
\(\varphi_i\) with every \(R_l\) replaced by \(\varphi_l^{(r)}\). By
\cref{thm:p03-w03-p03-059-pos-basic}, item~2, and induction on \(r\),
\(\varphi_i^{(r)}\) is a first-order formula defining the \(i\)-th
component of the \(r\)-th stage of \(F_{\mathcal S}\) in every structure.
On a one-element structure the set \(B\) of
\cref{def:p03-w03-p03-059-simultaneous-lfp} has \(m\) elements, so by
\cref{thm:p03-w03-p03-059-closure-ordinal} the closure ordinal is at most
\(m\) and \([\mathbf{lfp}\,R_j:\mathcal S](\bar t)\) is equivalent there to
the first-order formula \(\varphi_j^{(m)}(\bar t)\).

\emph{Item~2, conclusion.} Let \(\theta\in\mathrm{Pos}[\sigma]\). By steps
A and B it is equivalent to a formula \(\theta_B\) built by \(\land,\lor,
\exists,\forall\) from literals and top-level atoms of parameter-free
first-order systems. Let \(\theta_1\) arise from \(\theta_B\) by replacing
each atom by its first-order equivalent from step C2, and \(\theta_2\) by
replacing each atom by its \(\mathrm{LFP}_1\) equivalent from step C1. With
fresh \(u,v\), the formula
\[
 \bigl(\forall u\forall v\,u=v\land\theta_1\bigr)\lor
 \bigl(\exists u\exists v\,\neg u=v\land\theta_2\bigr)
\]
is equivalent to \(\theta\): on a one-element structure the first disjunct
reduces to \(\theta_1\) and the second is false, and on a larger structure
the first is false and the second reduces to \(\theta_2\). It is a formula
of \(\mathrm{LFP}_1\), being built by \(\land,\lor,\exists,\forall\) from
first-order formulas and least fixed points of parameter-free first-order
formulas positive in the fixed-point variable.

\emph{Item~3.} An LFP formula \(\chi\) of vocabulary \(\tau\) is equivalent
to \(\chi^+\), and \(\neg\chi\) to \(\chi^-\), by item~1; by item~2 both are
equivalent to formulas of \(\mathrm{LFP}_1\). Every formula of
\(\mathrm{LFP}_1\) is an LFP formula, since its fixed points are formed
from first-order formulas positive in the fixed-point variable and LFP is
closed under \(\land,\lor,\exists,\forall\). So if \(\psi\in\mathrm{LFP}_1\)
defines a query, \(\neg\psi\) is an LFP formula and is equivalent to a
formula of \(\mathrm{LFP}_1\), which defines the complementary query. A
greatest fixed point \([\mathbf{gfp}\,R\bar x\,.\,\varphi](\bar t)\) is an
LFP formula, so it too has an \(\mathrm{LFP}_1\) equivalent, in which
negation occurs only inside first-order subformulas.
\end{proof}
\end{annexitem}

\begin{annexitem}{annex:p03-w03-p03-059-pfp-dichotomy}{thm:p03-w03-p03-059-pfp-dichotomy}{The two possible behaviours of finite stages}
\begin{proof}
There are exactly \(N=2^{n^k}\) subsets of \(A^k\), and each stage
determines the next, since \(R^{m+1}=F_\psi(R^m)\).

\emph{Case 1: some stage is a fixed point.} Let \(m\) be least with
\(R^{m+1}=R^m\). Then \(R^{m'}=R^m\) for all \(m'\geq m\) by induction on
\(m'\). Suppose \(R^i=R^j\) for some \(i<j\leq m\). Then the sequence from
\(i\) on repeats with period \(j-i\): \(R^{i+r}=R^{j+r}\) for all
\(r\geq0\), by induction on \(r\). Choose \(r\) with \(i+r\geq m\) and
\(r\) a multiple of \(j-i\); then \(R^i=R^{i+r}=R^m\), and so
\(R^{i+1}=F_\psi(R^i)=F_\psi(R^m)=R^m=R^i\). Thus \(i<m\) is a fixed-point
stage, contradicting the choice of \(m\). So \(R^0,\ldots,R^m\) are
pairwise distinct subsets of \(A^k\), whence \(m+1\leq N\), that is,
\(m<N\). Moreover \(R^N=R^m\) is a fixed point.

\emph{Case 2: no stage is a fixed point.} The \(N+1\) stages
\(R^0,\ldots,R^N\) lie in a set of size \(N\), so \(R^i=R^j\) for some
\(i<j\leq N\), and as in case~1 the sequence is periodic from \(i\) with
period \(j-i\). The least period is not one, since a period one would make
\(R^i\) a fixed point.

The two cases exclude each other, which proves the dichotomy. For the
consequence: in case~1, \(R^N=R^m\) is the partial fixed point and
\(F_\psi(R^N)=R^N\); in case~2, the partial fixed point is \(\varnothing\)
by definition and \(F_\psi(R^N)\neq R^N\).
\end{proof}
\end{annexitem}

\begin{annexitem}{annex:p03-w03-p03-059-lfp-in-pfp}{thm:p03-w03-p03-059-lfp-in-pfp}{LFP is contained in PFP on finite structures}
\begin{proof}
\emph{The semantic step.} Let \(F\) be a monotone operator on \(A^k\), \(A\)
finite. By \cref{thm:p03-w03-p03-059-finite-stages} its finite stages,
which are the stages \(R^m\) of \cref{def:p03-w03-p03-059-pfp}, satisfy
\(R^m=R^{m+1}=\operatorname{lfp}(F)\) for some \(m\leq n^k\). So the
sequence reaches a fixed point, and the partial fixed point of \(F\) is
\(\operatorname{lfp}(F)\). The partial fixed point depends only on the
operator, not on the formula that defines it.

\emph{The displayed equivalence.} If \(R\) is positive in \(\psi\), then
\(F_\psi\) is monotone on every structure
(\cref{thm:p03-w03-p03-059-positive-monotone}), so its partial fixed point
is its least fixed point, and the two formulas hold of the same tuples.

\emph{Translation of LFP into PFP.} Define \(\chi^*\) for LFP formulas
\(\chi\) by recursion: atoms are unchanged; \(*\) commutes with the
connectives and quantifiers; and
\[
\begin{aligned}
 [\mathbf{lfp}\,R\bar x\,.\,\varphi](\bar t)^*&\defeq
   [\mathbf{pfp}\,R\bar x\,.\,\varphi^*](\bar t),\\
 [\mathbf{gfp}\,R\bar x\,.\,\varphi](\bar t)^*&\defeq{}\\
   &\neg[\mathbf{pfp}\,R\bar x\,.\,\neg\varphi^*[R/\neg R]](\bar t).
\end{aligned}
\]
We
show by induction on \(\chi\) that \(\chi\) and \(\chi^*\) are satisfied by
the same assignments in every finite structure, under every interpretation
of their free relation variables. Only the fixed-point steps need comment.
For \([\mathbf{lfp}\,R\bar x\,.\,\varphi](\bar t)\), the induction
hypothesis applied to \(\varphi\), for each interpretation of \(R\), gives
\(F_{\varphi^*}=F_\varphi\); this operator is monotone, so by the semantic
step its partial fixed point is \(\operatorname{lfp}(F_\varphi)\). For
\([\mathbf{gfp}\,R\bar x\,.\,\varphi](\bar t)\), the substitution claim in
the proof of \cref{thm:p03-w03-p03-059-gfp-dual} holds for PFP formulas
as well, since for a \(\mathbf{pfp}\) subformula it again shows that two
operators coincide, and equal operators have equal stages. Together with
the induction hypothesis \(F_{\varphi^*}=F_\varphi\) it gives, for every
\(X\),
\[
 F_{\neg\varphi^*[R/\neg R]}(X)=\overline{F_{\varphi^*}(\overline X)}
 =\overline{F_\varphi(\overline X)}=F_\varphi^d(X).
\]
The operator \(F_\varphi^d\) is monotone, so by the semantic step its
partial fixed point is \(\operatorname{lfp}(F_\varphi^d)=
\overline{\operatorname{gfp}(F_\varphi)}\)
(\cref{thm:p03-w03-p03-059-stages-lfp}, item~3). Hence
\(\neg[\mathbf{pfp}\,R\bar x\,.\,\neg\varphi^*[R/\neg R]](\bar t)\) holds
exactly when \(\bar t^{\mathfrak A}\in\operatorname{gfp}(F_\varphi)\). Hence every LFP
formula has an equivalent PFP formula on finite structures.
\end{proof}
\end{annexitem}

\begin{annexitem}{annex:p03-w03-p03-059-counter}{thm:p03-w03-p03-059-counter}{The counter reaches its fixed point at stage \(2^n-1\)}
\begin{proof}
List \(A\) in increasing order as \(e_0<e_1<\cdots<e_{n-1}\) and put
\(v(S)\defeq\sum_{e_j\in S}2^j\) for \(S\subseteq A\). Binary expansion
makes \(v\) a bijection from \(\mathcal P(A)\) onto
\(\{0,1,\ldots,2^n-1\}\), with \(v(A)=2^n-1\).

\emph{The top.} For \(S=A\) the third disjunct \(\forall y\,Ry\) is true at
every \(x\), so \(F_\psi(A)=A\).

\emph{The increment.} Let \(S\neq A\); then the third disjunct is false at
every \(x\). Let \(j\) be least with \(e_j\notin S\), so
\(e_0,\ldots,e_{j-1}\in S\). We show
\[
 F_\psi(S)=\bigl(S\setminus\{e_0,\ldots,e_{j-1}\}\bigr)\cup\{e_j\},
\]
by cases on \(x=e_i\).
\begin{itemize}
 \item \(i<j\): \(x\in S\) and every \(y<x\) is in \(S\), so the first
       disjunct is false; the second is false since \(x\in S\). So
       \(x\notin F_\psi(S)\).
 \item \(i=j\): \(x\notin S\) and every \(y<x\) is in \(S\), so the second
       disjunct holds and \(x\in F_\psi(S)\).
 \item \(i>j\) and \(x\in S\): \(y=e_j<x\) has \(y\notin S\), so the first
       disjunct holds and \(x\in F_\psi(S)\).
 \item \(i>j\) and \(x\notin S\): the first disjunct needs \(x\in S\); the
       second needs every \(y<x\) in \(S\), which fails at \(e_j\). So
       \(x\notin F_\psi(S)\).
\end{itemize}
Thus \(v(F_\psi(S))=v(S)-\sum_{i<j}2^i+2^j=v(S)-(2^j-1)+2^j=v(S)+1\).

\emph{The stages.} By induction on \(m\), \(v(R^m)=m\) for
\(0\leq m\leq2^n-1\): \(v(\varnothing)=0\), and if \(v(R^m)=m<2^n-1\) then
\(R^m\neq A\) and \(v(R^{m+1})=m+1\) by the increment. Since \(v\) is
injective, \(R^0,\ldots,R^{2^n-1}\) are pairwise distinct and
\(R^{2^n-1}=v^{-1}(2^n-1)=A\). For \(m<2^n-1\) we have
\(R^{m+1}\neq R^m\), so no earlier stage is a fixed point, while
\(R^{2^n}=F_\psi(A)=A=R^{2^n-1}\). The partial fixed point is therefore
\(A\), first reached at stage \(2^n-1\).
\end{proof}
\end{annexitem}

\begin{annexitem}{annex:p03-w03-p03-059-pfp-pspace}{thm:p03-w03-p03-059-pfp-pspace}{PFP formulas are evaluated in polynomial space}
\begin{proof}
Fix the sentence \(\psi\). On input a word \(u\), first test whether \(u\)
is a code and compute \(n=|A|\), in polynomial time by property~(a); if
\(u\) is not a code, reject. Elements are identified with
\(0,\ldots,n-1\) through the order of the code. By property~(c),
\(n\leq|u|\), so a bound polynomial in \(n\) is polynomial in \(|u|\).

We describe a recursive procedure \(\mathrm{Eval}(\chi,s,I)\) for
subformulas \(\chi\) of \(\psi\), where \(s\) assigns elements to the free
first-order variables of \(\chi\) and \(I\) assigns to each free relation
variable \(P\) of \(\chi\) of arity \(r\) a table of \(n^r\) bits stored on
the work tape. We show by induction on \(\chi\) that there is a polynomial
\(p_\chi\), depending only on \(\chi\), such that the procedure returns
the truth value of \(\chi\) under \(s\) and \(I\) using at most
\(p_\chi(n)\) work cells beyond the storage of \(s\) and \(I\), the input
being read-only.
\begin{itemize}
 \item An atom \(Q\bar y\) with \(Q\in\tau\) is decided by property~(a),
       in polynomial time and hence in polynomial space. An equality is
       decided by comparing two numbers below \(n\). An atom \(P\bar y\)
       with \(P\) a relation variable is decided by reading the bit of the
       table \(I(P)\) at the rank of \(s(\bar y)\) in the lexicographic
       order, computed with \(O(r\log n)\) cells.
 \item \(\neg\chi\), \(\chi_1\land\chi_2\), \(\chi_1\lor\chi_2\): evaluate
       the subformulas one after the other, reusing the space;
       \(p\defeq\max(p_{\chi_1},p_{\chi_2})+O(1)\).
 \item \(\exists x\chi\), \(\forall x\chi\): run a counter \(a\) through
       \(0,\ldots,n-1\) and evaluate \(\chi\) with \(x\mapsto a\), reusing
       the space; \(p\defeq p_\chi+O(\log n)\).
 \item \([\mathbf{pfp}\,R\bar x\,.\,\chi](\bar t)\) with \(R\) of arity
       \(k\): reserve two tables \(\mathit{Cur}\) and \(\mathit{Nxt}\) of
       \(n^k\) bits and a binary counter able to count to
       \(N=2^{n^k}\), which needs \(n^k+1\) bits. Set \(\mathit{Cur}\)
       to all zeros. Repeat \(N\) times: for each \(\bar a\in A^k\) in
       lexicographic order, set the \(\bar a\)-bit of \(\mathit{Nxt}\) to
       \(\mathrm{Eval}(\chi,s[\bar x\mapsto\bar a],I[R\mapsto\mathit{Cur}])\);
       then copy \(\mathit{Nxt}\) to \(\mathit{Cur}\). After this loop
       \(\mathit{Cur}=R^N\). Compute \(\mathit{Nxt}=F_\chi(\mathit{Cur})\)
       once more and compare. If they agree, return the bit of
       \(\mathit{Cur}\) at \(\bar t\); otherwise return false. By
       \cref{thm:p03-w03-p03-059-pfp-dichotomy} this is the truth value of
       the formula. The space used is
       \(p\defeq p_\chi+3n^k+O(k\log n)\).
\end{itemize}
Since \(\psi\) is fixed, it has finitely many subformulas, and each
\(p_\chi\) is a polynomial. The total work space is \(p_\psi(n)\) plus the
polynomial space of the initial test, which is polynomial in \(|u|\). The
running time may be exponential, because of the counters, but only space is
bounded. Hence the set of codes of finite models of \(\psi\) is in
\(\mathsf{PSPACE}\).

For a formula \(\psi(\bar x)\) with free first-order variables
\(x_1,\ldots,x_r\), the input is a code \(u\) together with \(r\) numbers in
binary. After the test of property~(a), reject unless each number is below
\(n\); this needs \(O(r\log|u|)\) further cells beyond the input. Then call
\(\mathrm{Eval}(\psi,s,\varnothing)\) with \(s(x_i)\) the \(i\)-th number.
The induction above did not use that \(\psi\) is a sentence, so the work
space is again \(p_\psi(n)\) plus the polynomial space of the tests, and
\(n\leq|u|\) by property~(c).
\end{proof}
\end{annexitem}

\begin{annexitem}{annex:p03-w03-p03-059-pfp-captures}{thm:p03-w03-p03-059-pfp-captures}{Abiteboul, Vianu and Vardi: PFP captures PSPACE on ordered structures}
\begin{proof}
If \(\mathcal K\) is the class of ordered models of a PFP sentence of
vocabulary \(\tau\cup\{<\}\), then \(\mathcal K\) is recognizable in
\(\mathsf{PSPACE}\) by \cref{thm:p03-w03-p03-059-pfp-pspace}, applied to
the vocabulary \(\tau\cup\{<\}\), whose atoms \(x<y\) are decided by
comparing ranks. For the converse, let \(\mathcal K\) be
closed under isomorphism with its set of codes in \(\mathsf{PSPACE}\).

\emph{Step 1: a single-tape machine.} By \cref{def:p04-resource-classes}
there is a deterministic machine \(M_0\), with a read-only input tape and a
work tape, that decides the set of codes of \(\mathcal K\) using at most
\(s_0(\ell)\) work cells on inputs of length \(\ell\), for a polynomial
\(s_0\), and that halts on every input. We turn it into a deterministic
single-tape machine \(M\) with state set \(Q\), tape alphabet
\(\Sigma\supseteq\Sigma_0\cup\{\sqcup\}\) (\(\sqcup\) the blank), start
state \(q_0\), and two halting states \(q_{\mathrm{acc}},q_{\mathrm{rej}}\).
The input occupies cells \(1,\ldots,\ell\), the head starts on cell~1, and
cell~1 is the left end: a move left from cell~1 leaves the head there.
\(M\) first rewrites each input cell as a two-track symbol (input symbol,
work symbol \(\sqcup\)), with marks for the two simulated heads on cell~1.
To simulate one step of \(M_0\), \(M\) sweeps from cell~1 rightward to the
first cell still holding the plain blank \(\sqcup\), recording in its
state the two marked symbols; it then computes the transition of \(M_0\)
and sweeps back, rewriting the marked work symbol and moving each mark by
at most one cell. When the work mark moves onto the plain blank cell, that
cell is converted to a two-track cell. When \(M_0\) halts, \(M\) enters
\(q_{\mathrm{acc}}\) or \(q_{\mathrm{rej}}\) accordingly. \(M\) uses at
most \(\max(\ell,s_0(\ell))+1\) cells, accepts exactly when \(M_0\) does,
and halts on every input. We may take \(s_0\) nondecreasing. Since \(\ell\leq p(n)\) by
property~(ii) of \cref{def:p03-w03-p03-059-good-encoding}, the polynomial
\(s(n)\defeq p(n)+s_0(p(n))+1\) satisfies \(\ell\leq s(n)\), and on the
code of an \(n\)-element structure \(M\) never visits a cell beyond
\(s(n)\).

Next modify the halting behaviour. An accepting configuration (state
\(q_{\mathrm{acc}}\)) is made its own successor. A new state
\(q_{\mathrm{rej}}'\) is added: a configuration in state
\(q_{\mathrm{rej}}\) passes to the same configuration in state
\(q_{\mathrm{rej}}'\), and back again. Write \(\mathrm{Next}(C)\) for the
successor configuration. Along the run of \(M\) on any input, a
configuration \(C\) with \(\mathrm{Next}(C)=C\) is accepting: a
non-halting \(C\) with \(\mathrm{Next}(C)=C\) would make the run stay at
\(C\) forever, contradicting that \(M\) halts; and rejecting configurations
change state.

\emph{Step 2: configurations as words.} Choose \(k\) with
\(n^k\geq s(n)+3\) for every \(n\geq2\); this is possible because
\(s(n)+3\leq c\,n^d\) for all \(n\geq2\) and suitable constants \(c,d\),
and then any \(k\) with \(2^{k-d}\geq c\) works. Put \(L\defeq n^k-1\). A
configuration with state \(q\), tape contents \(w_1w_2\cdots\) and head on
cell \(i\) is represented by the word \(c_0c_1\cdots c_L\) over
\[
 \Gamma\defeq\Sigma\cup(Q\times\Sigma)\cup\{\#\},
\]
where \(c_0=c_L=\#\), \(c_i=(q,w_i)\), and \(c_j=w_j\) for the other
\(1\leq j\leq L-1\). Since \(s(n)\leq n^k-3=L-2\), the head and every
non-blank cell lie in \(1,\ldots,L-2\).

Because the head moves at most one cell per step, the symbol at position
\(i\in\{1,\ldots,L-1\}\) of \(\mathrm{Next}(C)\) is a function of
\(c_{i-1}c_ic_{i+1}\). Define \(f:\Gamma^3\to\Gamma\setminus\{\#\}\) as
follows, where \(\delta(q,a)=(q',a',D)\) is the transition of \(M\) for a
non-halting state \(q\), with \(D\in\{\mathrm L,\mathrm R\}\); in cases 2
and 3 the state \(q\) is non-halting:
\begin{enumerate}
 \item if \(\beta=(q,a)\): if \(q=q_{\mathrm{acc}}\), then
       \(f(\alpha,\beta,\gamma)=(q_{\mathrm{acc}},a)\); if
       \(q=q_{\mathrm{rej}}\) or \(q=q_{\mathrm{rej}}'\), then the value is
       \((q_{\mathrm{rej}}',a)\) or \((q_{\mathrm{rej}},a)\)
       respectively; otherwise the value is \((q',a')\) when
       \(D=\mathrm L\) and \(\alpha=\#\) (the head is on cell~1 and
       stays), and \(a'\) in all other cases;
 \item if \(\beta\in\Sigma\) and \(\alpha=(q,a)\) with
       \(\delta(q,a)=(q',a',\mathrm R)\), then the value is \((q',\beta)\);
 \item if \(\beta\in\Sigma\) and \(\gamma=(q,a)\) with
       \(\delta(q,a)=(q',a',\mathrm L)\), then the value is \((q',\beta)\);
 \item if \(\beta\in\Sigma\) and neither of the two previous cases
       applies, the value is \(\beta\);
 \item on all remaining triples (\(\beta=\#\)) the value is \(\sqcup\).
\end{enumerate}
In a configuration only one cell carries a state, so cases 2 and 3 never
both apply. For every configuration \(C=c_0\cdots c_L\) on the run,
\(\mathrm{Next}(C)=c_0'\cdots c_L'\) with \(c_0'=c_L'=\#\) and
\(c_i'=f(c_{i-1},c_i,c_{i+1})\) for \(1\leq i\leq L-1\): the head cell is
rewritten (case~1), the cell the head moves to receives the new state
(cases 2 and 3; a move left from cell~1 is the exception in case~1), and
every other interior cell is unchanged (case~4).

\emph{Step 3: formulas.} Assume \(n\geq2\). In \((\mathfrak A,<)\), the
lexicographic order on \(k\)-tuples, the formulas \(\mathrm{First}(\bar y)\)
and \(\mathrm{Last}(\bar y)\) for the least and greatest \(k\)-tuple, and
the successor relation \(\bar x+1=\bar y\) on \(k\)-tuples are first-order
definable from \(<\); the rank of a \(k\)-tuple in that order is its
position \(0,\ldots,L\). Let \(\beta_\sigma(\bar x)\), \(\sigma\in\Sigma_0\),
be the formulas of property~(iii) for this \(k\). The input configuration
\(C_0\) has \(\#\) at \(0\) and \(L\), \((q_0,u_0)\) at \(1\) (with
\(u_0\) the first symbol of the code, or \(\sqcup\) if the code is empty),
the code symbol \(u_{i-1}\) at each position \(i\) with
\(2\leq i\leq\ell\), and \(\sqcup\) elsewhere. With
\(\mathrm{Code}(\bar x)\defeq\bigvee_{\sigma\in\Sigma_0}\beta_\sigma(\bar x)\),
the formulas
\[
\begin{aligned}
 \mathrm{In}_\#(\bar y)&\defeq\mathrm{First}(\bar y)\lor\mathrm{Last}(\bar y),\\
 \mathrm{In}_{(q_0,\sigma)}(\bar y)&\defeq\exists\bar x\bigl(\mathrm{First}(\bar x)
   \land\bar x+1=\bar y\land\beta_\sigma(\bar x)\bigr)\quad(\sigma\in\Sigma_0),\\
 \mathrm{In}_{(q_0,\sqcup)}(\bar y)&\defeq\exists\bar x\bigl(\mathrm{First}(\bar x)
   \land\bar x+1=\bar y\land\neg\mathrm{Code}(\bar x)\bigr),\\
 \mathrm{In}_{\sigma}(\bar y)&\defeq\exists\bar x\bigl(\neg\mathrm{First}(\bar x)
   \land\bar x+1=\bar y\land\beta_\sigma(\bar x)\bigr)\quad(\sigma\in\Sigma_0),\\
 \mathrm{In}_{\sqcup}(\bar y)&\defeq\neg\mathrm{Last}(\bar y)\land
   \exists\bar x\bigl(\neg\mathrm{First}(\bar x)\land\bar x+1=\bar y
   \land\neg\mathrm{Code}(\bar x)\bigr),
\end{aligned}
\]
and \(\mathrm{In}_\gamma\defeq(\bar y\neq\bar y)\) for every other
\(\gamma\in\Gamma\), define \(C_0\): \(\mathrm{In}_\gamma(\bar b)\) holds
exactly when \(C_0\) has \(\gamma\) at the position of \(\bar b\). (Because
\(\ell\leq s(n)\leq L-2\), the code has no symbol at the ranks \(L-2\) and
\(L-1\); so for positions \(L-1\) and \(L\) no \(\beta_\sigma\) applies,
position \(L-1\) receives \(\sqcup\), and only \(\mathrm{In}_\#\) holds at
\(L\).
If \(\sqcup\in\Sigma_0\), the formulas for \(\sigma=\sqcup\) and the blank
formula are merged by disjunction.)

The configuration is stored in one relation. Choose \(h\) with
\(2^h\geq|\Gamma|\) and an injective map \(\gamma\mapsto\ulcorner
\gamma\urcorner\in\{0,1\}^h\). As \(n\geq2\), the least and greatest
elements differ. Let \(\mathrm{Is}_\gamma(\bar u)\), for an \(h\)-tuple
\(\bar u\), say that \(u_i\) is the least element where the \(i\)-th bit of
\(\ulcorner\gamma\urcorner\) is \(0\) and the greatest element where it is
\(1\). For a relation variable \(T\) of arity \(h+k\) write
\[
 C_\gamma\bar y\defeq\exists\bar u\bigl(\mathrm{Is}_\gamma(\bar u)\land
 T\bar u\bar y\bigr),\qquad
 \mathrm{Empty}\defeq\forall\bar u\,\forall\bar y\,\neg T\bar u\bar y,
\]
and put
\[
 \psi(T,\bar u,\bar y)\defeq\bigvee_{\gamma\in\Gamma}
 \bigl(\mathrm{Is}_\gamma(\bar u)\land\theta_\gamma(T,\bar y)\bigr),
\]
where \(\theta_\#\defeq\mathrm{First}(\bar y)\lor\mathrm{Last}(\bar y)\)
and, for \(\gamma\neq\#\),
\[
\begin{aligned}
 \theta_\gamma\defeq{}&\bigl(\mathrm{Empty}\land\mathrm{In}_\gamma(\bar y)\bigr)
 \lor{}\\
 &\exists\bar x\,\exists\bar z\Bigl(\bar x+1=\bar y\land\bar y+1=\bar z
 \land{}\\
 &\qquad\bigvee_{f(\alpha,\beta,\varepsilon)=\gamma}
 \bigl(C_\alpha\bar x\land C_\beta\bar y\land C_\varepsilon\bar z\bigr)\Bigr).
\end{aligned}
\]

\emph{Step 4: the stages.} For a configuration \(C\) put
\(E(C)\defeq\{(\ulcorner\gamma\urcorner,\bar b):\text{\(C\) has \(\gamma\)
at the position of \(\bar b\)}\}\), with the bits read as least and greatest
elements. \(E\) is injective, and \(E(C)\neq\varnothing\). Let
\(C_0,C_1,\ldots\) be the run of \(M\) on \(\mathit{code}(\mathfrak A,<)\),
with \(C_{t+1}=\mathrm{Next}(C_t)\). We claim that the stages of
\(F_\psi\) satisfy \(T^{t+1}=E(C_t)\) for all \(t\geq0\).

For \(t=0\), \(T^0=\varnothing\), so \(\mathrm{Empty}\) holds and every
\(C_\alpha\) is empty. Then \(\theta_\#\) holds at the two end positions,
and for \(\gamma\neq\#\) the second disjunct of \(\theta_\gamma\) fails, so
\(\theta_\gamma\) reduces to \(\mathrm{In}_\gamma\). Hence
\(T^1=E(C_0)\). Assume \(T^{t+1}=E(C_t)\). Then \(\mathrm{Empty}\) fails,
and \(C_\alpha\bar b\) holds exactly when \(C_t\) has \(\alpha\) at
\(\bar b\). At an interior position \(i\), exactly one triple
\((\alpha,\beta,\varepsilon)\) satisfies the three atoms, namely
\((c_{i-1},c_i,c_{i+1})\) of \(C_t\), so \(\theta_\gamma\) holds exactly
for \(\gamma=f(c_{i-1},c_i,c_{i+1})\), which never equals \(\#\). At the
end positions the successor conditions fail, so only \(\theta_\#\) holds.
By Step~2, \(T^{t+2}=E(\mathrm{Next}(C_t))=E(C_{t+1})\). No tuple
\(\bar u\) outside the codes \(\ulcorner\gamma\urcorner\) ever enters a
stage.

\emph{Step 5: acceptance.} \(T^0=\varnothing\neq T^1\). For \(t\geq0\),
\(T^{t+1}=T^{t+2}\) holds exactly when \(C_t=C_{t+1}\), by injectivity of
\(E\), that is, exactly when \(C_t\) is accepting (Step~1). If \(M\)
accepts, the stages reach the fixed point \(E(C_{\mathrm{acc}})\neq
\varnothing\). If \(M\) rejects, the run ends in the two-cycle between the
rejecting states and no stage is a fixed point, so the partial fixed point
is \(\varnothing\). Hence, for \(n\geq2\),
\[
 (\mathfrak A,<)\in\mathcal K\iff
 (\mathfrak A,<)\entails\exists\bar u\,\exists\bar y\,
 [\mathbf{pfp}\,T\bar u\bar y\,.\,\psi](\bar u,\bar y).
\]

\emph{Step 6: one-element structures.} An ordered \(\tau\)-structure with one
element \(a\) is determined up to isomorphism by which relations contain the
constant tuple \((a,\ldots,a)\). For each of the finitely many such types
let \(\delta_j\) be the conjunction of \(\forall x\forall y\,x=y\) with
\(\forall x\,Qx\cdots x\) or \(\forall x\,\neg Qx\cdots x\) for each
\(Q\in\tau\). The sentence
\[
 \Bigl(\bigvee_{j:\ \text{type \(j\) is in \(\mathcal K\)}}\delta_j\Bigr)\lor
 \Bigl(\exists x\exists y\,x\neq y\land\exists\bar u\,\exists\bar y\,
 [\mathbf{pfp}\,T\bar u\bar y\,.\,\psi](\bar u,\bar y)\Bigr)
\]
is a PFP sentence of vocabulary \(\tau\cup\{<\}\), and by Step~5 and the
choice of the \(\delta_j\) its finite ordered models are exactly the members
of \(\mathcal K\). (The empty structure is not a structure here, since
universes are nonempty.)
\end{proof}
\end{annexitem}

\begin{annexitem}{annex:p03-w03-p03-059-ptime-pspace}{thm:p03-w03-p03-059-ptime-pspace}{\(\mathsf P=\mathsf{PSPACE}\) as a question about fixed points}
\begin{proof}
Fix an encoding with the properties (i)--(iv) and (a)--(c) of
\cref{def:p03-w03-p03-059-good-encoding}. By (a) and (b), any two such
encodings are translated into each other in polynomial time (decode the
size and the atoms, then encode), so membership of the set of codes of a
class in \(\mathsf P\) or in \(\mathsf{PSPACE}\) does not depend on which
of them is used. The capture of \(\mathsf P\) by LFP on ordered finite
structures \cite[Theorem~3.3.16, p.~157]{Gradel2007FMTDescriptiveChapter}
(\cref{thm:p04-immerman-vardi}; proved in Immerman's syntax as
\cref{thm:p09-immerman1986-capture}, whose binder
\((\bar t:\mu R(\bar x))\varphi\) is the formula
\([\mathbf{lfp}\,R\bar x\,.\,\varphi](\bar t)\)) says that an
isomorphism-closed class of ordered finite structures is recognizable in
\(\mathsf P\) exactly when it is defined by an LFP sentence.

\emph{(1) implies (2).} Every LFP sentence is equivalent on finite
structures to a PFP sentence by \cref{thm:p03-w03-p03-059-lfp-in-pfp}.
Conversely, if a class of ordered finite structures is defined by a PFP
sentence, it is recognizable in \(\mathsf{PSPACE}\) by
\cref{thm:p03-w03-p03-059-pfp-pspace}, hence in \(\mathsf P\) by (1), hence
defined by an LFP sentence by the capture of \(\mathsf P\).

\emph{(2) implies (1).} \(\mathsf P\subseteq\mathsf{PSPACE}\) because a
machine cannot use more cells than steps. Let \(\mathcal L\in\mathsf{PSPACE}\)
be a language over an alphabet \(\{a_1,\ldots,a_r\}\). Take the vocabulary
\(\tau_W=\{P_1,\ldots,P_r\}\) of unary symbols, and to a nonempty word
\(w=w_0\cdots w_{m-1}\) assign the word structure \(B(w)\) with universe
\(\{0,\ldots,m-1\}\), its natural order, and
\(P_i=\{j:w_j=a_i\}\) \cite[p.~137]{Gradel2007FMTDescriptiveChapter}. Let
\(\mathcal K\) be the class of ordered \(\tau_W\)-structures isomorphic to
\(B(w)\) for some nonempty \(w\in\mathcal L\). Given a word \(v\), decide
whether \(v\) is the code of a member of \(\mathcal K\) as follows: check
that \(v\) is a code and compute \(n\) (property~(a)); read off from the
atoms \(P_i(j)\) whether the structure is a word structure, that is,
whether every element lies in exactly one \(P_i\), and if so read off the
word \(w\) of length \(n\); then run the space-bounded decider for
\(\mathcal L\) on \(w\). This uses space polynomial in \(n\leq|v|\)
(property~(c)), so \(\mathcal K\) is recognizable in \(\mathsf{PSPACE}\).
By \cref{thm:p03-w03-p03-059-pfp-captures}, \(\mathcal K\) is defined by a
PFP sentence; by (2) it is defined by an LFP sentence; by the capture of
\(\mathsf P\), its set of codes is in \(\mathsf P\). Finally
\(\mathcal L\in\mathsf P\): on input \(w\), answer the fixed membership
of the empty word if \(w\) is empty, and otherwise compute the code of
\(B(w)\) in polynomial time (property~(b)) and test it for membership in
the set of codes of \(\mathcal K\). Hence
\(\mathsf{PSPACE}\subseteq\mathsf P\).
\end{proof}
\end{annexitem}
